% ============================================================
%  quiz.tex — Robot World: Meet the Machines That Help Us
%  End-of-book quiz (6 parts + answer key)
% ============================================================

\namedchapter{Robot World Quiz}
\addcontentsline{toc}{chapter}{Quiz}

Test what you have learned! Try each part, then check your answers at the end.
You can do this alone, with a partner, or as a team challenge.

\sectionrule

% --- Part 1: True or False ---
\section*{Part 1: True or False}

Write T (True) or F (False) next to each statement.

\begin{enumerate}
  \item A toaster is a robot. \hfill T \,/\, F
  \item Every robot follows the sense\,--\,think\,--\,act loop. \hfill T \,/\, F
  \item Lidar uses laser beams to measure distance. \hfill T \,/\, F
  \item All robots look like humans. \hfill T \,/\, F
  \item Cobots are designed to work safely beside people. \hfill T \,/\, F
  \item Mars rovers can be driven in real time from Earth. \hfill T \,/\, F
  \item Exoskeletons are wearable robots that help people move. \hfill T \,/\, F
  \item Swarm robots need one powerful leader robot. \hfill T \,/\, F
  \item A robot vacuum uses sensors to avoid falling down stairs. \hfill T \,/\, F
  \item The word ``robot'' comes from a Japanese word meaning ``helper.'' \hfill T \,/\, F
\end{enumerate}

\sectionrule

% --- Part 2: Multiple Choice ---
\section*{Part 2: Multiple Choice}

Circle the best answer.

\begin{enumerate}
  \item What three things does every robot do?\\
    a) Eat, sleep, work \quad b) Sense, think, act \quad c) Fly, drive, swim

  \item Which sensor uses sound waves to detect objects?\\
    a) Camera \quad b) Lidar \quad c) Ultrasonic sensor

  \item What does the controller do?\\
    a) Provides power \quad b) Runs the programme \quad c) Connects to Wi-Fi

  \item What is a cobot?\\
    a) A robot that works alone \quad b) A robot that works with humans \quad c) A robot pet

  \item What was the name of the first aircraft to fly on Mars?\\
    a) Spot \quad b) Ingenuity \quad c) Canadarm2

  \item What does sensor fusion mean?\\
    a) Sensors melting together \quad b) Combining data from many sensors \quad c) A sensor breaking

  \item Which of these is a soft robot inspired by?\\
    a) A brick wall \quad b) An octopus \quad c) A steel beam

  \item What is machine learning?\\
    a) A robot reading a textbook \quad b) A computer improving through experience \quad c) A human teaching a class
\end{enumerate}

\sectionrule

% --- Part 3: Match the Robot ---
\section*{Part 3: Match the Robot to Its Job}

Draw a line from each robot type to its job.

\begin{center}
\begin{tabularx}{0.85\textwidth}{|X|X|}
  \hline
  \textbf{Robot} & \textbf{Job} \\
  \hline
  Robot vacuum & Performs surgery \\
  \hline
  da Vinci system & Explores Mars \\
  \hline
  Perseverance & Cleans floors \\
  \hline
  AIBO & Welds car parts \\
  \hline
  Factory arm & Companion pet \\
  \hline
  Canadarm2 & Delivers packages in a warehouse \\
  \hline
  Warehouse robot & Catches spacecraft on the ISS \\
  \hline
\end{tabularx}
\end{center}

\sectionrule

% --- Part 4: Fill in the Blanks ---
\section*{Part 4: Fill in the Blanks}

Use words from the box to complete each sentence.

\begin{center}
\begin{tcolorbox}[colback=LightGray,colframe=RoboGreen!40,arc=4pt,boxrule=0.6pt,
  left=8pt,right=8pt,top=6pt,bottom=6pt]
  \centering
  chassis \quad sensor \quad algorithm \quad vehicle \quad bio \quad battery \quad programme \quad autonomous
\end{tcolorbox}
\end{center}

\begin{enumerate}
  \item A robot's body frame is called its \underline{\hspace{3cm}}.
  \item \underline{\hspace{3cm}} fusion means combining data from many sensors.
  \item An algorithm is a step-by-step method written into a \underline{\hspace{3cm}}.
  \item ROV stands for Remotely Operated \underline{\hspace{3cm}}.
  \item Robots inspired by animals are called \underline{\hspace{3cm}}-inspired.
  \item A robot that acts without human control is called \underline{\hspace{3cm}}.
  \item A rechargeable \underline{\hspace{3cm}} stores energy for mobile robots.
  \item A clear set of step-by-step instructions is called an \underline{\hspace{3cm}}.
\end{enumerate}

\sectionrule

% --- Part 5: Short Answer ---
\section*{Part 5: Short Answer}

Write 1--2 sentences for each question.

\begin{enumerate}
  \item Name two types of gripper a robot might use and explain when each is best.
  \item Why do Mars rovers need to be partly autonomous?
  \item Give one reason why companion robots help elderly people.
  \item What is the difference between an ROV and an AUV?
  \item Why must robot designers think about privacy and safety?
  \item How does a cobot know to stop if it bumps into a person?
\end{enumerate}

\sectionrule

% --- Part 6: Creative Challenge ---
\section*{Part 6: Creative Challenge}

Design a robot that solves a problem in your school or community.

\begin{itemize}
  \item Draw your robot and label its sensors, actuators, and brain.
  \item Write a three-step algorithm for its main task.
  \item Explain what energy source it uses and why.
  \item Give it a name and write one paragraph explaining why it would help.
  \item Bonus: What safety features would you include?
\end{itemize}

\sectionrule

% --- Answer Key ---
\section*{Answer Key}

\textbf{Part 1 --- True or False:}\\
1.\,F \quad 2.\,T \quad 3.\,T \quad 4.\,F \quad 5.\,T \quad 6.\,F \quad 7.\,T \quad 8.\,F \quad 9.\,T \quad 10.\,F (Czech, not Japanese)

\vspace{6pt}
\textbf{Part 2 --- Multiple Choice:}\\
1.\,b \quad 2.\,c \quad 3.\,b \quad 4.\,b \quad 5.\,b \quad 6.\,b \quad 7.\,b \quad 8.\,b

\vspace{6pt}
\textbf{Part 3 --- Match:}\\
Robot vacuum \textrightarrow{} Cleans floors;
da Vinci \textrightarrow{} Performs surgery;
Perseverance \textrightarrow{} Explores Mars;
AIBO \textrightarrow{} Companion pet;
Factory arm \textrightarrow{} Welds car parts;
Canadarm2 \textrightarrow{} Catches spacecraft on the ISS;
Warehouse robot \textrightarrow{} Delivers packages in a warehouse.

\vspace{6pt}
\textbf{Part 4 --- Fill in the Blanks:}\\
1.\,chassis \quad 2.\,Sensor \quad 3.\,programme \quad 4.\,Vehicle \quad 5.\,bio \quad 6.\,autonomous \quad 7.\,battery \quad 8.\,algorithm

\vspace{6pt}
\textbf{Part 5 --- Sample Answers:}
\begin{enumerate}
  \item Pincer gripper (good for solid objects like bolts); suction cup (good for flat surfaces like glass panels). Also accept: magnetic, soft, multi-finger.
  \item Radio signals take up to 20 minutes to travel between Earth and Mars, so the rover cannot wait for instructions --- it must make some decisions on its own.
  \item They reduce loneliness, provide comfort, and can lower stress by giving residents something to interact with.
  \item An ROV is connected to a ship by cable and controlled by a human pilot; an AUV navigates on its own with no cable.
  \item A robot with cameras or microphones could collect private information; engineers must ensure robots only gather what is needed, keep data safe, and include stop buttons.
  \item Cobots have force sensors that detect unexpected contact and instantly stop all movement to prevent injury.
\end{enumerate}
